\documentclass[exercice,a4]{thedocument}%
\startdatakeys
    \source{laclassecooperativedestef.com}
    \BO{2026}
    \cycle{3}
    \competencesDuSocle{RE4,CO2}
    \objectifsApprentissage{C3C1ab,C3C1aj,C3C2ab,C3C1ae,C3C2an,C3C1ag}
    \motsCles{rectangle, milieu, perpendicularite, parallelisme}
\enddatakeys
\begin{document}%
\noindent\textbf{Programme de construction 25}\par\smallskip\noindent
\begin{itemize}
    \item Tracer un rectangle $\mathspoint{ABCD}$ de longueur
    $\mathspoint{AB}=6\text{ cm}$ et de largeur $\mathspoint{BC}=2\text{ cm}$.
    \item Placer $\mathspoint{E}$, milieu de $[\mathspoint{AD}]$, et
    $\mathspoint{F}$, milieu de $[\mathspoint{BC}]$.
    \item À droite de $[\mathspoint{BC}]$, tracer le segment
    $[\mathspoint{FJ}]$ tel que $\mathspoint{FJ}=2\text{ cm}$,
    perpendiculaire à $[\mathspoint{BC}]$.
    \item À gauche de $[\mathspoint{AD}]$, tracer le segment
    $[\mathspoint{IE}]$ tel que $\mathspoint{IE}=2\text{ cm}$,
    perpendiculaire à $[\mathspoint{AD}]$.
    \item Placer le point $\mathspoint{G}$ sur $[\mathspoint{DC}]$ à
    $2\text{ cm}$ de $\mathspoint{D}$.
    \item Placer le point $\mathspoint{H}$ sur $[\mathspoint{DC}]$ à
    $2\text{ cm}$ de $\mathspoint{C}$.
    \item En dessous de $[\mathspoint{DC}]$, tracer $[\mathspoint{HL}]$ tel
    que $\mathspoint{HL}=1\text{ cm}$, perpendiculaire à $[\mathspoint{DC}]$.
    \item En dessous de $[\mathspoint{DC}]$, tracer $[\mathspoint{GK}]$ tel
    que $\mathspoint{GK}=1\text{ cm}$, perpendiculaire à $[\mathspoint{DC}]$.
    \item Tracer le segment $[\mathspoint{NK}]$ tel que
    $\mathspoint{NK}=4\text{ cm}$, perpendiculaire à $[\mathspoint{GK}]$.
    \item Tracer le segment $[\mathspoint{LM}]$ tel que
    $\mathspoint{LM}=4\text{ cm}$, perpendiculaire à $[\mathspoint{HL}]$.
    \item Tracer la droite $(\mathspoint{d})$ passant par
    $\mathspoint{G}$ et $\mathspoint{H}$.
    \item Placer $\mathspoint{O}$ sur $(\mathspoint{d})$ sachant que
    $\mathspoint{OD}=1{,}5\text{ cm}$.
    \item Placer $\mathspoint{P}$ sur $(\mathspoint{d})$ sachant que
    $\mathspoint{CP}=1{,}5\text{ cm}$.
    \item Tracer les segments $[\mathspoint{JP}]$, $[\mathspoint{PM}]$,
    $[\mathspoint{LC}]$, $[\mathspoint{DK}]$, $[\mathspoint{IO}]$ et
    $[\mathspoint{ON}]$.
    \item Colorier les triangles en bleu foncé et les autres figures en
    bleu clair.
\end{itemize}
\end{document}%
